\section{Discussion}
\label{sec:discussion}

\paragraph{Limitations} All confirmatory evidence comes from one environment that we designed to defeat shortcuts. It is small (24 tasks) and nearly deterministic, and its development and evaluation episodes share a generator. The candidate causes are given, and \textit{other} is only a residual; generating new hypotheses is outside this work. Our verifier is strict: a verdict must cite an observation carrying the cause's own signature, so a correct verdict reached by elimination counts as a failure. This rule applies to every configuration since the initial study, makes our completion rates conservative, and is the entire cost of the controller-side stop. We do not evaluate evidence manipulation or prompt injection through log content. An attacker who can make a probe return the outcome expected under a benign cause can still mislead the ledger, and in our environment the planner's reading of raw text never changes a score. We tested two model families, five sizes, one temperature, and one serving stack, and public data could not provide a real-data confirmation (\S\ref{sec:realdata}). Finally, we made errors of our own: we misreported one interval, lost the journals of one study, found a ranking leak into a control arm and a table asymmetry that favoured the baseline, and corrected the ground truth of a descriptive metric twice. Appendix~\ref{app:errata} documents each of them; none affects a confirmed primary endpoint.

\paragraph{Where the Model Is Still Needed} When observations arrive as discrete outcomes and the tables come from development data, the procedure does most of the work: the controller alone completes 0.917 of cases, and small models inside the controller do no better. Two roles remain for the model. The first is judging when evidence suffices. Strong models that decide when to stop reach 1.000 because they wait for a confirming signature, a judgment that could also be written into the controller as a cost-aware stopping rule. The second is producing the discrete observations in the first place. A real probe returns log lines, not the label \textit{high entropy}, and turning one into the other requires reading. Our environment performs that translation itself, so our results do not measure it. The OTRF case in \S\ref{sec:realdata} shows what it involves: recognizing that a task's creator is a person rather than a machine account, or that a logon came over the network.

\paragraph{Design Alternatives} Each design choice replaced an alternative that we implemented and measured. Publishing the ranking and letting the planner choose changes completion by at most $+0.06$ once the controller selects, so the controller executes its own choice. A budget-aware two-step lookahead selector beat EIG/cost by $+0.17$ to $+0.21$ on the task family we developed it on, but on a held-out family (3{,}888 episodes) it matched or trailed EIG/cost, by $-0.125$ for the 7B planner, and on sec-triage it completed 0.889, 0.972, and 0.958 against 0.958, 1.000, and 1.000 for the three Qwen models; we kept the greedy rule. Showing the planner a posterior above 0.999 and stating that no legal probe is left does not make Llama-3.1-8B conclude, whereas a stopping rule enforced by the controller does. Without the finish guard, completion changes by at most $+0.04$, but the guard ensures that every accepted verdict cites current supporting evidence, so we keep it.

\paragraph{Potential Applications} \hesp{} can (i) run triage end to end on a local model and hand each verdict, with its cited evidence and journal, to an analyst, or (ii) act as a copilot that proposes the next probe and flags when the evidence meets the acceptance rule. A deployment cannot treat ``let the controller choose'' as free: for Llama-3.1-8B the missing piece was the stopping rule, for Qwen2.5-7B both stopping and probe selection, and for the 32B and 72B models taking the choice away was harmful. The prediction tables are the main deployment cost. In our environment they came from development episodes of the same generator; in practice they would come from resolved tickets, expert estimates, or an LLM prior corrected by data, and which of these is accurate enough remains open.

\paragraph{Runtime and Cost} End-to-end latency is dominated by LLM inference; ledger updates and selection take negligible time next to a model call. The running example took 5.1\,s for three planner calls on an 8B model, and the stopping study's 1{,}800 episodes over five models ran in 28 minutes on one GPU with 32 concurrent episodes. Controller selection also reduces probe cost: with a controller stop, it halves Llama-3.1-8B's cost from 4.78 to 2.35 units per episode, and the LLM-free controller has the lowest cost of any configuration.
